\documentclass[12pt]{article}
\usepackage[utf8]{inputenc}
\usepackage[T1]{fontenc}
\usepackage{hyperref}
\title{Tiwi Language Resources for K-12 Curriculum Development}
\author{Compiled by Claude Sonnet 5 (claude-sonnet-5) for Andr\'as Kornai}
\date{}
\begin{document}
\maketitle
\section{Introduction}
Tiwi (ISO 639-3 tiw; Glottocode tiwi1244) is a non-Pama-Nyungan language of the Tiwi Islands (Bathurst and Melville), northern Australia, with no known close relatives. Traditional Tiwi is mainly spoken by older people, and modern Tiwi has absorbed English influence. One search was run; current school programmes were not verified.
\section{Primary Educational Resources}
An encyclopedia entry (1986-era data) says the Nguiu bilingual literature centre on Bathurst Island produced Tiwi texts for the local primary school, while Milikapiti and Parlingimpi schools taught in English: \url{https://www.encyclopedia.com/humanities/encyclopedias-almanacs-transcripts-and-maps/tiwi}. This may be outdated.
\section{Linguistic and Reference Materials}
C.\,R. Osborne (1974), \emph{The Tiwi Language} (grammar, myths, dictionary): \url{https://glottolog.org/resource/reference/id/697300}. Jennifer Lee's 1993 Tiwi-English dictionary (Nguiu Nginingawila Literature Production Centre) was mentioned in results; no direct URL seen. AIATSIS holds Lee's data books, including dictionary manuscript and children's stories in Tiwi: \url{https://aiatsis.gov.au/sites/default/files/catalogue_resources/MS2381.PDF}. Language record: \url{https://collection.aiatsis.gov.au/austlang/language/n20}.
\section{Religious and Community Resources}
Cultural background: \url{https://www.everyculture.com/Oceania/Tiwi-Orientation.html}. No scripture or community-organisation pages were found.
\section{Digital and Audio Resources}
None found.
\section{Curriculum Development Context}
Archival materials (Osborne, Lee) exist and a literature production centre historically supported Tiwi in school. Current status of bilingual education should be confirmed with Tiwi community bodies and the Northern Territory education department.
\end{document}